%======================================================================
\section{A purity budget at the rank rate}\label{sec:purity-budget}
%======================================================================

At the rank rate a device cannot buy memory with purity: the purity it can spend is linear in its memory. This is the last tool for the lower
bound at the rank rate in Theorem~B. We sharpen the budget $14Q+3$ of~\cite{DouglasMghirbiB} by combining its rank count with a one-sided
second-moment bound.

\begin{proposition}[Purity budget]\label{prop:purity-budget}
Let $\Phi$ be a channel on $M_d$, and suppose that some pure state $\psi_*$ of the input and a reference has output
$\omega=(\Phi\otimes\mathrm{id})(\psi_*)$ flat of rank $r$. Every causal service of $\Phi$ for $n$ rounds with error $\epsilon\le1/16$ and exchange
$J\le\frac n2\log r$ satisfies
\[
P+C\ <\ 3Q+6 .
\]
\end{proposition}

Every gadget channel has such a probe with $r=r_*$ (Section~\ref{sec:channel-setting}). The proof uses the variance of the surprisal of a
state with small deficit, which is~\cite[Lemma~5.9]{DouglasMghirbiB}; we include the short argument.

\begin{lemma}[Surprisal variance]\label{lem:surprisal}
Let $\beta$ be a state on $q\ge1$ qubits with spectrum $(\beta_i)$ and deficit $D=q-S(\beta)$. The surprisal $-\log\beta_i$, drawn with probability
$\beta_i$, has variance at most $(q+2/\ln2)\,D$.
\end{lemma}

\begin{proof}
Put $\varphi(x)=x\ln x-x+1\ge0$. We claim $x(\ln x)^2\le(2+\ln_+x)\varphi(x)$ for $x\ge0$. For $x=e^t$ with $t\ge0$ the difference of the two sides
is $f(t)=e^t(t-2)+t+2$, with $f(0)=f'(0)=0$ and $f''(t)=te^t\ge0$. For $t\le0$ it is $g(t)=2(te^t-e^t+1)-t^2e^t$, with $g(0)=0$ and
$g'(t)=-t^2e^t\le0$; and $x=0$ is clear. Now let $Q_0=2^q$ and $x_i=Q_0\beta_i\le Q_0$, the index running over all $Q_0$ eigenvalues. Then
$\frac1{Q_0}\sum_ix_i=1$ and $D\ln2=\sum_i\beta_i\ln(Q_0\beta_i)=\frac1{Q_0}\sum_i\varphi(x_i)$. The variance is at most
\[
\sum_i\beta_i\bigl(\log(Q_0\beta_i)\bigr)^2=\frac1{Q_0(\ln2)^2}\sum_ix_i(\ln x_i)^2\le\frac{2+q\ln2}{(\ln2)^2}\,D\ln2=\Bigl(q+\frac2{\ln2}\Bigr)D .
\]
\end{proof}

\begin{proof}[Proof of Proposition~\ref{prop:purity-budget}]
Put $b=P+C$ and suppose, for a contradiction, that $b\ge3Q+6$.

\emph{Preloading.} The imports are fixed states, independent of everything else, and the exported qubits are never used again. So the observer's
final state does not change if every import is prepared at the start in a register of its own, and each replacement swaps it with the exported
qubits, which then stay in the device. The device becomes a register of $L=Q+J$ qubits in the product state
$\sigma_0=\beta_0\otimes\bigotimes_f\omega_f$ of the initial memory state and the imported states, and it applies a fixed unitary in each round.

\emph{The surprisal.} Let $(\sigma_i)$ be the spectrum of $\sigma_0$ and let $Z=L+\log\sigma_i$, with $i$ drawn with probability $\sigma_i$. Then
$\mathbb EZ=L-S(\sigma_0)=b$. Also $Z$ is the sum of the independent variables $Z_f=q_f+\log\beta^{(f)}_i$ attached to the factors of $\sigma_0$,
where $q_f$ is the number of qubits of the factor and $\beta^{(f)}$ its spectrum. Every factor has at most $Q$ qubits, since an import replaces
exported memory qubits, and the deficits $b_f=\mathbb EZ_f$ of the factors add up to $b$. By Lemma~\ref{lem:surprisal},
\[
\mathrm{Var}\,Z=\sum_f\mathrm{Var}\,Z_f\le\Bigl(Q+\frac2{\ln2}\Bigr)b .
\]

\emph{A tail bound from the rank.} Let the observer feed the input halves of $n$ copies of $\psi_*$ and keep the references. Its final state is
$\tau_O=\sum_i\sigma_i\tau_i$, where $\tau_i$ is its final state when the register starts in the eigenvector of $\sigma_i$. Each $\tau_i$ is a
marginal of a pure state of the observer's system and the $L$-qubit register, so $\rank\tau_i\le2^L$. Fix a real number $k$ and let $I_k$ be the
set of indices with $\sigma_i\ge2^{k-L}$, so that $|I_k|\le2^{L-k}$. Let $E$ be the support projection of $\sum_{i\in I_k}\tau_i$. Then
$\rank E\le2^{2L-k}$ and $\Tr(E\tau_O)\ge\sum_{i\in I_k}\sigma_i=\Pr(Z\ge k)$. The ideal final state $\omega^{\otimes n}$ is flat of rank $r^n$, so
$\Tr(E\omega^{\otimes n})\le2^{2L-k}r^{-n}\le2^{2Q-k}$, using $2J\le n\log r$. Since the error is at most $\epsilon$,
\begin{equation}\label{eq:rank-tail}
\Pr(Z\ge k)\ \le\ \epsilon+2^{2Q-k}\qquad(k\in\mathbb R).
\end{equation}

\emph{The contradiction.} Put $u=b-2Q-3$, so that $u\ge b/3+1>0$. By the one-sided Chebyshev inequality,
$\Pr(Z\ge b-u)\ge u^2/(\mathrm{Var}\,Z+u^2)$, while~\eqref{eq:rank-tail} with $k=b-u=2Q+3$ gives $\Pr(Z\ge b-u)\le\frac1{16}+\frac18=\frac3{16}$. Hence
$13u^2\le3\,\mathrm{Var}\,Z$, and with $u>b/3$,
\[
\frac{13}9\,b^2<3\Bigl(Q+\frac2{\ln2}\Bigr)b,\qquad\text{that is}\qquad b<\frac{27}{13}\Bigl(Q+\frac2{\ln2}\Bigr)<2.08\,Q+5.993 ,
\]
which contradicts $b\ge3Q+6$.
\end{proof}
